\chapter{Protein C and Protein S}

\section*{What Is This Test?}
Protein C and protein S are vitamin K-dependent plasma proteins that function together as a natural anticoagulant system. Protein C is a serine protease zymogen synthesized in the liver that, when activated by the thrombin-thrombomodulin complex on the surface of intact endothelial cells, becomes activated protein C (APC). APC, together with its cofactor protein S, inactivates activated coagulation factors Va (cofactor in the prothrombinase complex) and VIIIa (cofactor in the tenase complex) by proteolytic cleavage at specific arginine residues. This anticoagulant pathway serves as a critical feedback inhibition mechanism that localizes thrombin generation to sites of vascular injury while preventing excessive clot propagation into normal vessels \cite{reitsma2012mechanistic}. Protein C also enhances fibrinolysis by inhibiting plasminogen activator inhibitor-1 (PAI-1). Deficiency of protein C or protein S, whether inherited (autosomal dominant) or acquired, results in an increased risk of venous thromboembolism (VTE), with affected individuals having a 2--11 fold increased risk of deep vein thrombosis, pulmonary embolism, cerebral venous sinus thrombosis, and, in the most severe cases, warfarin-induced skin necrosis and neonatal purpura fulminans \cite{dahlback2008advances}. Protein C and S levels are measured by functional (activity) and immunologic (antigen) assays \cite{tierney2023current}.

\section*{Alternative Names}
Protein C Activity; Protein C Antigen; Protein S Activity; Protein S Antigen; Free Protein S; Total Protein S; Natural Anticoagulant Testing; Thrombophilia Panel Components

\section*{Principle}
\textbf{Protein C} functional activity is measured by a chromogenic assay. Patient plasma is incubated with a specific protein C activator (derived from snake venom -- Protac, from Agkistrodon contortrix) that converts protein C to activated protein C (APC). A synthetic chromogenic substrate specific for APC is added, and APC cleaves the substrate, releasing para-nitroaniline (pNA). The color change is measured spectrophotometrically at 405 nm and is proportional to the protein C concentration. Clot-based protein C assays using aPTT- or PT-based systems are also available but are less commonly used. Protein C antigen (immunologic) is measured by ELISA or latex immunoturbidimetric assay. \textbf{Protein S} is more complex to measure because approximately 60\% of protein S in plasma is bound to C4b-binding protein (C4bBP), and only the free fraction (approximately 40\%) functions as the APC cofactor. Protein S functional activity is measured by a clot-based assay: patient plasma is added to protein S-depleted plasma with APC and factor Va, and the prolongation of the clotting time is proportional to protein S cofactor activity. Free protein S antigen is measured by latex immunoturbidimetric assay or ELISA using a monoclonal antibody that specifically recognizes free protein S (not bound to C4bBP), or by polyethylene glycol (PEG) precipitation of the C4bBP-bound fraction followed by measurement of the remaining (free) fraction. Total protein S antigen measures both free and bound forms. The free protein S antigen and protein S activity assays are the clinically relevant tests; total protein S antigen measurement adds little \cite{tietz2012textbook}.

\section*{History}
Protein C was discovered by Johan Stenflo in 1976 in Lund, Sweden, and named ``protein C'' because it was the third protein peak (C) eluting from a chromatographic column. Its role as an anticoagulant was elucidated in the late 1970s and early 1980s by Charles Esmon and colleagues, who demonstrated that protein C was activated by the thrombin-thrombomodulin complex and that activated protein C inactivated factors Va and VIIIa \cite{esmon2003protein}. Protein S was discovered by Richard DiScipio in 1977 in Seattle and named after Seattle (S). The clinical significance of protein C and S deficiencies was established in the early 1980s, when John H. Griffin and colleagues described protein C deficiency as a cause of familial thrombophilia, and Philip Comp and Charles Esmon described protein S deficiency and the phenomenon of warfarin-induced skin necrosis. The association between protein C deficiency and neonatal purpura fulminans was described by Branson in 1983. The landmark PROWESS trial (2001) demonstrated that recombinant human activated protein C (drotrecogin alfa, Xigris) reduced mortality in severe sepsis, leading to FDA approval. However, the PROWESS-SHOCK trial (2011) failed to confirm benefit, and the drug was withdrawn from the market in 2011.

\section*{Why This Test Is Ordered}
\begin{itemize}
  \item Diagnosis of hereditary protein C or protein S deficiency in patients with unexplained venous thromboembolism (DVT, PE), especially: young age at first thrombosis (<50 years), recurrent thrombosis, thrombosis at unusual sites (cerebral venous sinus, mesenteric vein, portal vein, hepatic vein), family history of VTE
  \item Evaluation of warfarin-induced skin necrosis risk -- protein C deficiency predisposes to this complication when warfarin is initiated without heparin bridging (protein C has a short half-life of approximately 6--8 hours and falls before the procoagulant factors, creating a transient hypercoagulable state)
  \item Investigation of neonatal purpura fulminans (homozygous or compound heterozygous protein C or S deficiency)
  \item Thrombophilia testing in recurrent pregnancy loss (second and third trimester losses, intrauterine fetal demise)
  \item Thrombophilia workup in premenopausal women with VTE considering hormonal contraception
  \item Screening of asymptomatic family members of patients with known protein C or S deficiency
  \item Assessment of liver synthetic function (protein C and S are synthesized in the liver; levels decrease in liver disease)
  \item Evaluation of DIC (consumption of protein C and S) and sepsis-associated purpura fulminans
  \item Investigation of coumarin-induced skin necrosis
\end{itemize}

\section*{Primary vs Secondary Test}
Protein C and protein S assays are secondary (specialized) thrombophilia tests ordered after VTE has occurred, in selected patients with specific clinical features suggesting inherited thrombophilia.

\section*{How This Test Is Performed}
A venous blood sample is collected in a light blue-top tube containing 3.2\% sodium citrate (9:1 ratio), filled exactly to the line \cite{fischbach2023manual}. The sample must be processed within 1 hour of collection. The plasma is separated by centrifugation at 1,500--2,500 $\times$ g for 15 minutes to obtain platelet-poor plasma. For functional assays: protein C chromogenic assay (incubate with Protac activator, add chromogenic substrate, measure absorbance at 405 nm); protein S clotting assay (mix patient plasma with protein S-depleted plasma, APC, factor Va, measure clotting time prolongation). For antigen assays: free protein S antigen by latex immunoturbidimetric assay or ELISA. If testing cannot be performed immediately, plasma must be frozen at -20\textdegree{}C. Venipuncture takes less than 5 minutes. Analysis takes 30--60 minutes.

\section*{What Preparation Is Needed}
No special preparation is required. Fasting is not necessary. The timing of testing is critical: protein C and S levels should NOT be measured during an acute thrombotic event (acute thrombosis consumption lowers levels, producing falsely low results); they should be measured 2--4 weeks after completion of anticoagulation therapy whenever possible. Warfarin therapy markedly decreases protein C and S levels (they are vitamin K-dependent) and results obtained during warfarin treatment are uninterpretable. Heparin and DOACs may also affect clot-based functional assays. Pregnancy decreases protein S levels physiologically (free protein S falls to 40--60\% of normal in the second and third trimesters) and protein S should not be assessed during pregnancy. Oral contraceptives and hormone replacement therapy decrease protein S levels. Liver disease, DIC, and recent surgery/trauma decrease protein C and S through consumption. All of these confounders must be accounted for when interpreting results.

\section*{Contraindications and Precautions}
No absolute contraindications. Critical considerations: the test MUST NOT be performed during warfarin therapy -- warfarin directly reduces protein C and S by inhibiting vitamin K-dependent gamma-carboxylation; protein C has the shortest half-life (6--8 hours) of the vitamin K-dependent proteins, followed by factor VII (4--6 hours), so protein C falls first during warfarin initiation, causing a transient hypercoagulable state. Testing during warfarin produces falsely low results and leads to misdiagnosis of deficiency. Tests should be performed ideally 2--4 weeks after completing anticoagulation. For patients who cannot safely stop warfarin, bridging to low-molecular-weight heparin (which does not affect protein C and S levels) is an option, but this is logistically complex and rarely done. Protein S levels are physiologically decreased during pregnancy -- up to 60\% reduction in free protein S -- and pregnancy should be excluded before testing. Protein S levels are also decreased by oral contraceptives and estrogen therapy. The sample for protein C and S testing must be double-centrifuged to ensure platelet-poor plasma (residual platelets release platelet factor 4, which neutralizes heparin and can affect clotting assays). Testing should be performed within 4 hours, or plasma stored at -20\textdegree{}C. Free protein S antigen is more clinically relevant than total protein S (bound to C4bBP is inactive).

\section*{Risks}
Minimal risks of venipuncture: mild pain, bruising, hematoma, lightheadedness, and rarely infection or nerve injury.

\section*{What Happens After the Test}
Results are available within 1--4 hours. If protein C or S deficiency is identified, the diagnosis must be confirmed by repeat testing at least 2--4 weeks after completion of anticoagulation and resolution of any acute thrombosis. Protein C activity <55--65\%, protein S activity <55--65\%, or free protein S antigen <55--65\% (age- and sex-adjusted) with normal PT and liver function tests suggests deficiency. If confirmed, family screening may be offered. The finding of protein C or S deficiency influences the duration of anticoagulation after VTE (many patients with deficiency are offered indefinite anticoagulation after a first provoked or unprovoked VTE) and may affect recommendations for thromboprophylaxis during high-risk situations (surgery, pregnancy, prolonged immobility). Asymptomatic individuals with protein C or S deficiency are generally not anticoagulated but should receive aggressive thromboprophylaxis during high-risk situations \cite{comp2021protein}.

\section*{Understanding Results}

\subsection*{Below Normal (Protein C or Protein S Deficiency)}
\begin{itemize}
  \item \textbf{Hereditary protein C deficiency:} Autosomal dominant. Prevalence 0.2--0.4\% in the general population; 3--5\% of patients with VTE. Heterozygous deficiency (activity 40--65\%): increased risk of VTE (7-fold relative risk), warfarin-induced skin necrosis. Homozygous deficiency (activity <1\%): presents in the neonatal period as purpura fulminans -- massive, rapidly progressive thrombosis of dermal microvasculature leading to skin necrosis, DIC, and death without prompt protein C replacement. Mutations in the PROC gene on chromosome 2 \cite{goldenberg2008protein}.
  \item \textbf{Hereditary protein S deficiency:} Autosomal dominant. Prevalence 0.1--0.3\%; 2--5\% of VTE patients. Three types: Type I (quantitative deficiency: low total and free protein S antigen), Type II (qualitative deficiency: normal antigen but low functional activity), Type III (low free protein S with normal total protein S, due to increased binding to C4bBP). Heterozygous: increased VTE risk (2--11 fold). Homozygous: neonatal purpura fulminans. Mutations in the PROS1 gene on chromosome 3 \cite{tenkate2008protein}.
  \item \textbf{Warfarin therapy (acquired):} Vitamin K antagonism reduces gamma-carboxylation of protein C and S, lowering functional levels. Protein C falls first due to its short half-life (6--8 hours), creating a transient prothrombotic state that causes warfarin-induced skin necrosis (typically on day 3--5 of warfarin initiation) in patients with underlying protein C deficiency. This is prevented by heparin bridging during warfarin initiation.
  \item \textbf{Vitamin K deficiency (acquired):} Dietary deficiency, malabsorption, antibiotic suppression of gut flora. Reduced protein C and S functional activity.
  \item \textbf{Liver disease:} Impaired hepatic synthesis. Protein C and S decrease proportional to synthetic dysfunction. Protein C is a marker of liver synthetic function.
  \item \textbf{DIC:} Widespread activation of coagulation consumes protein C and S.
  \item \textbf{Sepsis and severe infection:} Inflammatory cytokines (TNF-alpha, IL-1) downregulate thrombomodulin and endothelial protein C receptor (EPCR), impairing protein C activation. Sepsis-associated purpura fulminans, particularly in meningococcemia (Neisseria meningitidis), is associated with severe acquired protein C deficiency.
  \item \textbf{Pregnancy:} Physiologically decreased protein S (free protein S falls to 40--60\% of normal in second and third trimesters due to estrogen effects and increased C4bBP). Protein S should not be tested during pregnancy.
  \item \textbf{Oral contraceptives and estrogen therapy:} Decreased free protein S antigen and activity.
  \item \textbf{Nephrotic syndrome:} Urinary loss of protein S (smaller molecular weight than protein C), leading to selective protein S deficiency.
  \item \textbf{L-asparaginase chemotherapy:} Reduces hepatic synthesis of protein C and S.
  \item \textbf{Acute thrombosis:} Consumption during acute VTE can transiently lower levels. Testing should be delayed 2--4 weeks.
\end{itemize}

\subsection*{Above Normal (Elevated Protein C or Protein S)}
\begin{itemize}
  \item \textbf{Anabolic steroid use:} Elevated protein C and S levels.
  \item \textbf{Nephrotic syndrome (protein C):} Protein C (larger molecular weight than protein S, 62 kDa vs. 75 kDa, but retained due to size) may be elevated due to compensatory hepatic synthesis.
  \item \textbf{Advanced age:} Mild elevation of protein C and S with age, but clinical significance is unknown.
  \item \textbf{No known hypercoagulability risk:} Elevated protein C or S is not known to cause bleeding or other clinical consequences. No bleeding disorder is associated with elevated protein C or S.
\end{itemize}

\section*{Normal Results}
\begin{itemize}
  \item \textbf{Protein C activity (chromogenic):} Adults: 70--140\%
  \item \textbf{Protein C antigen (immunologic):} Adults: 70--140\%
  \item \textbf{Protein S activity (functional):} Adults: 60--140\% (males); 50--130\% (females) -- females have physiologically lower protein S
  \item \textbf{Free protein S antigen:} Adult males: 65--140\%; Adult females: 50--130\%; Adult females on oral contraceptives: 40--100\%
  \item \textbf{Total protein S antigen:} 60--140\%
  \item \textbf{Newborns:} Protein C and S are physiologically low at birth (protein C 20--40\%, protein S 20--50\%), reaching adult levels by 3--6 months (protein C) and 6--12 months (protein S)
  \item \textbf{Pregnancy:} Protein C is normal or mildly elevated; protein S decreases (free protein S 30--60\% of normal in second and third trimester)
  \item \textbf{Patients on warfarin:} Both protein C and S are decreased to approximately 40--60\% of normal at therapeutic INR -- results are invalid for diagnostic purposes
  \item \textbf{Diagnostic thresholds for hereditary deficiency:} Protein C activity <55--65\%; Protein S activity <55--65\%; Free protein S antigen <55--65\% (must be confirmed by repeat testing) \cite{wallach2022interpretation}
\end{itemize}

\section*{Follow-up Tests}
\begin{itemize}
  \item \textbf{Repeat protein C and S testing:} In 2--4 weeks (off warfarin and resolved acute thrombosis) to confirm deficiency
  \item \textbf{Activated protein C resistance (APC-R) assay:} To detect factor V Leiden, which is distinct from protein C deficiency but phenotypically results in resistance to APC
  \item \textbf{Factor V Leiden DNA mutation testing:} Factor V Leiden (R506Q) is the most common inherited thrombophilia (prevalence 3--7\% in Caucasians); APC resistance without protein C deficiency
  \item \textbf{Antithrombin activity:} The other major natural anticoagulant; antithrombin deficiency causes thrombophilia
  \item \textbf{Prothrombin G20210A mutation testing:} Second most common inherited thrombophilia after factor V Leiden
  \item \textbf{Antiphospholipid antibody testing:} Lupus anticoagulant, anticardiolipin antibodies, and anti-beta-2-glycoprotein I for acquired thrombophilia
  \item \textbf{Liver function tests and INR:} To rule out liver disease or warfarin effect as cause of low protein C/S
  \item \textbf{Serum protein electrophoresis (SPEP):} To rule out nephrotic syndrome (protein S loss) or monoclonal gammopathy
  \item \textbf{Genetic testing (PROC and PROS1 genes):} For definitive diagnosis of hereditary deficiency
  \item \textbf{DNA sequencing for PROC and PROS1 mutations:} Available in specialized laboratories; identifies specific mutations in families
  \item \textbf{Family screening:} Testing of first-degree relatives of patients with confirmed deficiency
\end{itemize}

\section*{References}
\bibliographystyle{unsrtnat}
\bibliography{references}

\clearpage
\section*{Regulatory Status and Authorization Date}
Regulatory status applies to the specific assay, instrument, manufacturer, and intended use rather than to the test name alone. No single approval date can be assigned to this test category. For a commercial assay, verify the current product in the relevant regulator's database: the U.S. Food and Drug Administration (FDA) clearance, approval, or authorization; India's Central Drugs Standard Control Organisation (CDSCO) licence or permission; the European Union In Vitro Diagnostic Regulation (EU IVDR) CE certificate issued through the applicable conformity-assessment route; or the United Kingdom Medicines and Healthcare products Regulatory Agency (MHRA) registration and UKCA/CE status. Laboratory-developed tests may not have an individual FDA, CDSCO, EU IVDR, or MHRA approval date.
\clearpage
